\renewcommand{\Anno}{2006}

\renewcommand{\NumeroProblema}{1}

\renewcommand{\Testo}{Determinare tutte le coppie $(x,y)$ di numeri \emph{interi} tali
	che
	$$x^4+3x^{2}y^{2}+9y^{4}=12^{2006}.$$}

\renewcommand{\Risposta}{$(0,2^{1003}\cdot 3^{501})$, 
$(0,-2^{1003}\cdot 3^{501})$.}

\renewcommand{\Aiutino}{È possibile che $x$ ed $y$ siano entrambi
dispari? E uno pari e uno dispari? Se sono entrambi pari si pone $x=2a$, $y=2b$ e dopo aver eliminato una fattore 16 ci si ritrova un'equazione analoga. Di questo passo spariscono tutti i fattori 2.

Passiamo ora ai fattori 3. Può $x$ non essere divisibile per 3? Se $x=3a$ si può eliminare un 9 ottenendo un'equazione simile a quella di partenza. Procedendo allo stasso modo prima o poi rimane a
secondo membro solo un 9. A questo punto è fatta.}

\renewcommand{\SoluzioneUfficiale}{\TitoloSoluzione Riscriviamo
l'equazione nella forma
$$x^4+3x^{2}y^{2}+9y^{4}=2^{4012}\cdot 3^{2006},$$
e trattiamo separatamente i fattori 2 ed i fattori 3.

È evidente che $x$ ed $y$ non possono essere entrambi dispari (il primo membro sarebbe somma di 3 numeri dispari, quindi dispari a sua volta), e non possono nemmeno essere uno pari ed uno dispari (anche in questo caso il primo membro sarebbe dispari, in quanto somma di 2 pari ed un dispari). Ne segue che $x$ ed $y$ sono entrambi pari. Possiamo pertanto porre $x=2a$, $y=2b$. Sostituendo nell'equazione otteniamo che
$$16a^4+3\cdot 16x^{2}y^{2}+9\cdot 16y^{4}=2^{4012}\cdot 3^{2006},$$
da cui, semplificando il 16, si ricava che
$$a^4+3x^{2}y^{2}+9y^{4}=2^{4012-4}\cdot 3^{2006}.$$

Abbiamo quindi ottenuto un'equazione analoga a quella di partenza, soltanto con 4 fattori 2 in meno al secondo membro.  Ripetendo lo stesso ragionamento $4012:4=1003$ volte possiamo eliminare tutti i fattori 2.  Più precisamente avremo che $x=2^{1003}\cdot m$, $y=2^{1003}\cdot n$, con $m$ ed $n$ numeri interi che soddisfano
$$m^4+3m^{2}n^{2}+9n^{4}=3^{2006}.$$

Passiamo ora ai fattori 3. È chiaro che $m$ deve essere multiplo di 3. Sostituendo $m=3h$ otteniamo che
$$3^4\cdot h^{4}+3^{3}m^{2}n^{2}+9n^{4}=3^{2006},$$
da, cui semplificando un 9,
$$9\cdot h^{4}+3m^{2}n^{2}+n^{4}=3^{2006-2}.$$

Ora è chiaro che $n$ deve essere multiplo di 3.  Sostituendo $n=3k$ e semplificando si ha che
$$h^4+3h^{2}k^{2}+9k^{4}=3^{2006-4},$$
che è ancora la stessa equazione, soltanto con 4 fattori 3 in meno al secondo membro.

Poiché $2006=501\cdot 4+2$, ripetendo 501 volte il procedimento rimarranno soltanto 2 fattori 3. Più precisamente avremo che $m=3^{501}\alpha$, $n=3^{501}\beta$, con $\alpha$ e $\beta$ numeri interi che soddisfano
$$\alpha^4+3\alpha^{2}\beta^{2}+9\beta^{4}=9.$$

È ora evidente che le uniche soluzioni di questa equazione sono le 2 coppie $(\alpha,\beta)=(0,\pm 1)$, da cui, ricordando le posizioni fatte, $(m,n)=(0,\pm 3^{501})$,
ed infine $(x,y)=(0,\pm 2^{1003}\cdot 3^{501})$.}

\renewcommand{\SoluzioniAlternative}{}

\renewcommand{\Osservazioni}{
}

\renewcommand{\Commenti}{
}